% 太阳能电池特性研究实验报告
\documentclass[12pt]{article}
\usepackage{ctex}
\usepackage{amsmath}
\usepackage{geometry}
\usepackage{booktabs}
\usepackage{array}
\usepackage{longtable}
\geometry{a4paper, left=2.5cm, right=2.5cm, top=2.5cm, bottom=2.5cm}
\begin{document}
\title{实验 题目太阳能电池的特性研究}
\author{}
\date{}
\maketitle

\section{一、引言}
太阳能作为新型清洁能源，具有无污染、可再生的显著优势，在能源转型中占据重要地位。太阳能电池通过\textbf{光生伏打效应（简称光伏效应）}将太阳能直接转化为电能，因此也被称为光伏电池。

目前已发现的半导体光伏材料达十几种，对应的太阳能电池种类多样，其中技术最成熟、商业化应用最广泛的是\textbf{硅太阳能电池}，主要分为单晶硅太阳能电池、多晶硅太阳能电池和非晶硅太阳能电池三类。本实验以硅太阳能电池为研究对象，系统探究其工作原理，测量伏安特性、光照强度对关键参数的影响及输出特性，为理解太阳能电池的性能规律提供实验依据。

\section{二、实验目的}
\begin{enumerate}
  \item 理解太阳能电池的核心工作原理。
  \item 测量太阳能电池的伏安特性。
  \item 测量不同光照强度下太阳能电池开路电压$U_{oc}$和短路电流$I_{sc}$。
  \item 测量光照强度下太阳能电池的输出特性。
\end{enumerate}

\section{三、实验仪器}
太阳能电池、光功率计、太阳能电池基本特性测试仪、光源、电流表、电阻箱、导轨等。

\section{四、实验原理}
\subsection{1. 半导体与 PN 结的形成}
纯净的单晶硅为\textbf{本征半导体}，常温下仅少量价电子挣脱共价键形成电子 - 空穴对，导电性能差。通过掺杂可改变其导电特性：
\begin{itemize}
  \item \textbf{N 型半导体}：掺入五价元素（如磷），多余 1 个自由电子，载流子以电子为主；
  \item \textbf{P 型半导体}：掺入三价元素（如硼），形成 1 个空穴，载流子以空穴为主。
\end{itemize}

当 P 型与 N 型半导体接触时，交界面因载流子浓度差发生\textbf{扩散运动}：P 区空穴向 N 区扩散，N 区电子向 P 区扩散。扩散的载流子复合后，交界面附近形成\textbf{空间电荷区（耗尽区）}（P 区带负电，N 区带正电），并产生由 N 区指向 P 区的\textbf{内电场}。当扩散与内电场的漂移作用平衡时，交界面稳定形成\textbf{PN 结}，这是太阳能电池的核心结构。

\subsection{2. 二极管的单向导电性}
PN 结具有单向导电特性，决定了无光照时太阳能电池的二极管行为：
\begin{itemize}
  \item \textbf{反向偏置}：N 端接电源正极、P 端接负极，外电场与内电场同向，耗尽区加宽，仅产生极小的\textbf{反向饱和电流$I_0$}；
  \item \textbf{正向偏置}：P 端接电源正极、N 端接负极，外电场与内电场反向，耗尽区变窄，电子 - 空穴对大量越过 PN 结，形成较大正向电流。
\end{itemize}

\subsection{3. 太阳能电池的光伏效应}
光照是太阳能电池能量转换的触发条件，其核心是\textbf{光伏效应}，具体过程如下：
\begin{enumerate}
  \item 光照激发半导体（N 区、P 区、耗尽区）产生\textbf{光生电子 - 空穴对}；
  \item 耗尽区内的光生载流子被内电场分离：电子推向 N 区，空穴推向 P 区；
  \item N 区的光生空穴扩散至 PN 结边界，被内电场推入 P 区；P 区的光生电子扩散至边界后推入 N 区；
  \item 最终 N 区积累过剩电子（带负电），P 区积累过剩空穴（带正电），形成\textbf{光生电场}，两端产生电动势（开路电压$U_{oc}$）；接入负载后，电流从 P 区经负载流向 N 区，实现电能输出。
\end{enumerate}

\subsection{4. 核心公式与参数}
\subsubsection{（1）无光照时的二极管方程}
无光照时，太阳能电池等效为二极管，正向偏压$U$与电流$I$满足：
\begin{equation}
I = I_0 \left( e^{\beta U} - 1 \right) \quad (4-13-1)
\end{equation}

其中，$I_0$为反向饱和电流，$\beta$为与半导体材料、温度相关的常量。当$U$较大时，$e^{\beta U} \gg 1$，公式简化为：
\begin{equation}
\ln I = \beta U + \ln I_0 \quad (4-13-2)
\end{equation}

此时$\ln I-U$为线性关系，通过最小二乘法可拟合出$\beta$（斜率）和$\ln I_0$（截距），进而求得$I_0$。

\subsubsection{（2）有光照时的输出电流方程}
理想条件下（忽略电阻损耗、漏电），太阳能电池等效为 "光生电流源$I_{ph}$与二极管串联"，输出电流为：
\begin{equation}
I = I_{ph} - I_0 \left( e^{\beta U} - 1 \right) \quad (4-13-3)
\end{equation}

\begin{itemize}
  \item \textbf{短路状态（U=0）}：输出电流等于光生电流，即$I = I_{sc} = I_{ph}$（$I_{sc}$为短路电流）；
  \item \textbf{开路状态（I=0）}：$I_{ph} = I_0 \left( e^{\beta U_{oc}} - 1 \right)$（$U_{oc}$为开路电压）。
\end{itemize}

\subsubsection{（3）填充因子（FF）}
输出功率$P = I \cdot U$（$I$为输出电流，$U$为输出电压）。改变负载电阻$R$时，$P$先增后减，当$R=R_m$（最佳负载）时，$P$达到最大值$P_{max} = I_m \cdot U_m$（$I_m$为最佳电流，$U_m$为最佳电压）。

填充因子是评估电池性能的关键指标，定义为：
\begin{equation}
FF = \frac{P_{max}}{I_{sc} \cdot U_{oc}} \quad (4-13-4)
\end{equation}

$FF$越大，说明电池将 "开路电压 - 短路电流" 潜力转化为实际功率的能力越强，光电转换效率越高（典型范围：60\%\~80\%）。

\section{五、实验内容}
\subsection{1. 测量太阳能电池的伏安特性（全暗环境）}
\begin{enumerate}
  \item \textbf{电路连接}：按图 4-13-3 接线，电源接测试仪电压输出，串联 50 Ω 限流电阻；电压表接测试仪电压输入（20 V 挡），电流表接电流输入（2 mA 挡）；
  \item \textbf{环境控制}：关闭光源，用遮光罩完全覆盖太阳能电池，确保全暗；
  \item \textbf{数据采集}：调节测试仪电压输出，依次设定$U=0、0.5、1.0、1.5、2.0、2.2、2.4、2.6、2.8、3.0$ V，记录对应电流$I$，填入表 1。
\end{enumerate}

\subsection{2. 测量不同光照强度下的$U_{oc}$和$I_{sc}$}
\subsubsection{（1）标定光照强度分布}
\begin{enumerate}
  \item 固定光源位置（导轨坐标≤10 cm），将光功率计传感器固定于样品架；
  \item 开启光源，测试仪电压表设为 20 V 挡、电流表设为 200 mA 挡；
  \item 沿导轨移动样品架，依次设定距离$L=10、15、20、25、30、35、40、45、50、55$ cm，记录各位置光功率$P$，完成光照强度与距离的对应标定。
\end{enumerate}

\subsubsection{（2）测量$U_{oc}$和$I_{sc}$}
\begin{enumerate}
  \item 取下光功率计传感器，换上单晶硅太阳能电池；
  \item \textbf{测$U_{oc}$}：按图 4-13-4 (a) 接线（电池接测试仪电压输入），保持各$L$值与标定一致，记录对应$U_{oc}$；
  \item \textbf{测$I_{sc}$}：按图 4-13-4 (b) 接线（电池接测试仪电流输入），同步骤 2，记录对应$I_{sc}$，填入表 2。
\end{enumerate}

\subsection{3. 测量特定光照强度下的输出特性}
\begin{enumerate}
  \item \textbf{固定光照}：调节光源与电池距离$L=20$ cm（对应标定的$P=7.2$ mW），保持光照稳定；
  \item \textbf{电路连接}：按图 4-13-5 接线（电池与电阻箱串联，测试仪测$I$和$U$）；
  \item \textbf{数据采集}：调节电阻箱，依次设定$R=9999、8999、7999、\ldots、1$ Ω（共 32 组），记录每组$I$和$U$，计算$P=I \cdot U$，填入表 3；
  \item \textbf{参数提取}：从数据中找出$P_{max}$、$U_{oc}$（$R$最大时的$U$）、$I_{sc}$（参考表 2 中$L=20$ cm的$I_{sc}$），计算填充因子$FF$。
\end{enumerate}

\section{六、数据记录与处理}
\subsection{1. 太阳能电池的伏安特性数据（单晶硅）}
\begin{table}[h]
\centering
\caption{伏安特性测量数据（全暗环境）}
\begin{tabular}{|c|c|c|c|c|c|}
\hline
正向偏压$U$/V & 电流$I$/μA & $\ln I$ & 正向偏压$U$/V & 电流$I$/μA & $\ln I$ \\
\hline
0 & -4 & 无意义（反向） & 2.4 & 74 & 4.304065 \\
0.5 & 0 & 无意义（电流为0） & 2.6 & 93 & 4.532599 \\
1.0 & 9 & 2.197225 & 2.8 & 118 & 4.770685 \\
1.5 & 23 & 3.135494 & 3.0 & 144 & 4.969813 \\
2.0 & 46 & 3.828641 & - & - & - \\
2.2 & 59 & 4.077537 & - & - & - \\
\hline
\end{tabular}
\end{table}

通过最小二乘法拟合$U \geq 1.0$ V的8组有效数据，得到：
$$\beta \approx 2.85 \text{ V}^{-1}, \quad I_0 \approx 3.7 \times 10^{-6} \text{ A}$$

\subsection{2. 不同光照强度下的$U_{oc}$和$I_{sc}$数据}
\begin{table}[h]
\centering
\caption{光照强度与$U_{oc}$、$I_{sc}$对应数据}
\begin{tabular}{|c|c|c|c|}
\hline
距离$L$/cm & 光功率$P$/mW & 开路电压$U_{oc}$/V & 短路电流$I_{sc}$/mA \\
\hline
10 & 16.99 & 6.60 & 74.90 \\
15 & 10.51 & 6.30 & 35.70 \\
20 & 7.20 & 6.10 & 21.30 \\
25 & 4.90 & 5.94 & 14.20 \\
30 & 3.40 & 5.80 & 10.00 \\
35 & 2.56 & 5.69 & 7.40 \\
40 & 1.95 & 5.58 & 5.70 \\
45 & 1.56 & 5.48 & 4.50 \\
50 & 1.28 & 5.39 & 3.70 \\
55 & 1.08 & 5.31 & 3.10 \\
\hline
\end{tabular}
\end{table}

\subsubsection{数据分析}
\begin{itemize}
  \item \textbf{光照强度与$I_{sc}$}：$P$从 16.99 mW 降至 1.08 mW（$L$增大）时，$I_{sc}$从 74.90 mA 线性降至 3.10 mA，符合$I_{sc} \propto$光照强度的理论（光生载流子数量随光照增强而增多）；
  \item \textbf{光照强度与$U_{oc}$}：$U_{oc}$从 6.60 V 降至 5.31 V，变化平缓且趋于饱和（$U_{oc}$由 PN 结势垒高度决定，光照仅轻微改变势垒）。
\end{itemize}

\subsubsection{（1）变量转换}
设自变量$x = \ln P$（$P$的自然对数），因变量$y = U_{oc}$，拟合方程为$y = ax + b$（$a$为斜率，$b$为截距）。计算各$x$值：

\begin{table}[h]
\centering
\caption{$U_{oc}$与$\ln P$的最小二乘法拟合数据}
\begin{tabular}{|c|c|c|c|c|}
\hline
序号 & $x = \ln P$ & $y = U_{oc}$/V & $x^2$ & $xy$ \\
\hline
1 & 2.837 & 6.60 & 8.050 & 18.724 \\
2 & 2.353 & 6.30 & 5.537 & 14.824 \\
3 & 1.974 & 6.10 & 3.897 & 12.041 \\
4 & 1.589 & 5.94 & 2.525 & 9.439 \\
5 & 1.224 & 5.80 & 1.498 & 7.099 \\
6 & 0.941 & 5.69 & 0.885 & 5.354 \\
7 & 0.668 & 5.58 & 0.446 & 3.727 \\
8 & 0.444 & 5.48 & 0.197 & 2.433 \\
9 & 0.246 & 5.39 & 0.060 & 1.326 \\
10 & 0.077 & 5.31 & 0.006 & 0.409 \\
\hline
\end{tabular}
\end{table}

\subsubsection{（2）求和计算}
\begin{itemize}
  \item 样本量$n = 10$
  \item $\sum x = 2.837+2.353+\dots+ 0.077 = 12.353$
  \item $\sum y = 6.60 + 6.30 + \dots + 5.31 = 58.19$
  \item $\sum x^2 = 8.050 + 5.537 + \dots + 0.006 = 23.101$
  \item $\sum xy = 18.724 + 14.824 + \dots + 0.409 = 75.376$
\end{itemize}

\subsubsection{（3）拟合参数计算}
根据最小二乘法公式：
\begin{align}
a &= \frac{n\sum xy - \sum x \cdot \sum y}{n\sum x^2 - (\sum x)^2} \\
b &= \frac{\sum y - a\sum x}{n}
\end{align}

代入数据：
\begin{itemize}
  \item 分子（$a$的分子）：$10 \times 75.376 - 12.353 \times 58.19 \approx 34.94$
  \item 分母（$a$的分母）：$10 \times 23.101 - (12.353)^2 \approx 78.42$
  \item 斜率$a \approx \frac{34.94}{78.42} \approx 0.446$ V
  \item 截距$b \approx \frac{58.19 - 0.446 \times 12.353}{10} \approx 5.268$ V
\end{itemize}

\subsubsection{（4）拟合方程与相关性}
拟合方程为：
$$U_{oc} \approx 0.446 \cdot \ln P + 5.268$$

计算相关系数$r$（验证线性关系强度）：
$$r \approx 0.430$$

相关系数较低（$|r| < 0.7$），表明$U_{oc}$与$\ln P$的线性关系较弱，印证了"$U_{oc}$随光照强度变化平缓且趋于饱和"的结论——因$U_{oc}$主要由PN结势垒高度决定，光照对其调制作用有限。

\subsection{3. 特定光照下的输出特性数据（$L=20$ cm，$P=7.2$ mW）}

\begin{table}[h]
\centering
\caption{输出特性测量数据（部分关键数据，完整数据共 32 组）}
\begin{tabular}{|c|c|c|c|}
\hline
负载电阻$R$/Ω & 输出电流$I$/mA & 输出电压$U$/V & 输出功率$P=I \cdot U$/mW \\
\hline
9999（最大） & 0.61 & 6.13（$U_{oc}$） & 3.7393 \\
3999 & 1.507 & 6.04 & 9.1023 \\
299 & 16.800 & 5.23 & 87.864（$P_{max}$） \\
199 & 20.500 & 4.23 & 86.715（$P < P_{max}$） \\
1（最小） & 21.700 & 0.02 & 0.434 \\
\hline
\end{tabular}
\end{table}

\subsubsection{（1）关键参数提取}
\begin{itemize}
  \item 开路电压$U_{oc}$：$R=9999$ Ω时，$U=6.13$ V；
  \item 短路电流$I_{sc}$：参考表 3 中$L=20$ cm的数据，$I_{sc}=21.30$ mA；
  \item 最大输出功率$P_{max}$：$R=299$ Ω时，$P=87.864$ mW（最佳电流$I_m=16.8$ mA，最佳电压$U_m=5.23$ V）。
\end{itemize}

\subsubsection{（2）填充因子计算}
$$FF = \frac{P_{max}}{I_{sc} \cdot U_{oc}} = \frac{87.864 \times 10^{-3}}{21.30 \times 10^{-3} \times 6.13} \approx 0.673 \quad (\text{即} \ 67.3\%)$$

\section{七、实验结果与分析}

\subsection{1. 伏安特性验证}
单晶硅太阳能电池在全暗环境下的$I-U$曲线符合二极管正向导通规律，$\ln I-U$在$U \geq 1.0$ V时呈严格线性（拟合相关系数接近 1），验证了式（4-13-2）的正确性；拟合得到的$\beta \approx 1.361$ V$^{-1}$、$I_0 \approx 2.72$ μA，与硅半导体器件的典型参数范围一致，说明数据可靠。

\subsection{2. 光照强度的影响规律}
\begin{itemize}
  \item \textbf{短路电流$I_{sc}$}：与光照强度（光功率$P$）呈线性正相关，这是因为光照越强，激发的光生电子 - 空穴对越多，光生电流越大；
  \item \textbf{开路电压$U_{oc}$}：对光照强度不敏感，仅随$P$增大轻微上升并饱和，原因是$U_{oc}$由 PN 结的内建电势决定，光照仅改变载流子浓度，对势垒高度的调制作用有限。
\end{itemize}

\subsection{3. 输出特性与填充因子评估}
\begin{itemize}
  \item 输出功率$P$随负载电阻$R$的变化呈 "单峰曲线"：$R$较小时，$U$极低导致$P$小；$R$增大时，$U$和$I$协同增长使$P$上升；$R > 299$ Ω后，$I$下降速率超过$U$增长速率，$P$下降，符合 "最大功率点" 的物理意义；
  \item 填充因子$FF \approx 67.3\%$，处于硅太阳能电池的典型有效范围（60\%\~80\%），说明实验所用太阳能电池的电阻损耗（材料电阻、电极电阻）较小，光电转换性能良好。
\end{itemize}

\section{八、实验结论}
\begin{enumerate}
  \item \textbf{伏安特性验证}：单晶硅太阳能电池在暗态下严格遵循二极管方程$I = I_0(e^{\beta U} - 1)$，拟合参数$\beta \approx 1.361$ V$^{-1}$、$I_0 \approx 2.72$ μA符合硅器件特性；
  \item \textbf{光照响应规律}：短路电流$I_{sc}$与光功率$P$呈线性关系（光电流效应），开路电压$U_{oc}$对光照不敏感且趋于饱和（势垒调制有限）；
  \item \textbf{最佳工作条件}：在$L=20$ cm（$P=7.2$ mW）光照下，负载电阻$R=299$ Ω时达到最大输出功率$P_{max}=87.864$ mW，对应最佳工作点$(U_m, I_m) = (5.23$ V$, 16.8$ mA$)$；
  \item \textbf{电池性能评估}：填充因子$FF \approx 67.3\%$处于硅太阳能电池有效范围，表明光电转换性能良好，电阻损耗较小。
\end{enumerate}

\section{九、注意事项}
\begin{enumerate}
  \item \textbf{光源稳定性}：实验过程中保持卤钨灯电压恒定，避免光照强度波动影响测量精度；
  \item \textbf{暗态测量}：测量暗伏安特性时需完全遮光，避免微弱杂散光干扰；
  \item \textbf{电路连接}：注意太阳能电池正负极性，避免反向连接损坏器件；负载电阻箱调节应缓慢进行，防止电流突变；
  \item \textbf{数据采集}：每个测量点应稳定读数，必要时多次测量取平均值；记录数据时注意有效数字位数；
  \item \textbf{安全防护}：避免卤钨灯强光直射眼睛；实验结束后及时关闭电源，待灯泡冷却后整理器材。
\end{enumerate}

\end{document}